\section{Background and Related Work}
\label{sec:background}

\subsection{Classical fingerprinting}

Traffic analysis without payload inspection has a long line of work built on
packet sizes, directions and inter-arrival times. CUMUL~\cite{panchenko2016website}
interpolates the cumulative signed-size curve to a fixed-length vector.
AppScanner~\cite{taylor2016appscanner} computes eighteen summary statistics over
the incoming, outgoing and combined size series and fits a random forest.
$k$-fingerprinting~\cite{hayes2016kfingerprinting} adds burst structure and
concentration features. All three predate the deep-learning literature in this
area by several years, and we find in \S\ref{sec:results} that AppScanner
remains within 0.02 macro-F1 of a modern neural model on one of our corpora.

\subsection{Neural classifiers}

Deep Packet~\cite{lotfollahi2020deeppacket} applies a 1D convolutional network
to the first 1500 payload bytes. FS-Net~\cite{liu2019fsnet} encodes discretised
packet lengths with a stacked bidirectional GRU. More recently, pre-trained
transformers such as ET-BERT~\cite{lin2022etbert} tokenise raw datagram bytes
and report macro-F1 above 98\% on several standard benchmarks.

\paragraph{Why we do not compare against ET-BERT.} Its input is the hexadecimal
byte string of the packet payload, tokenised as bigrams over a 65{,}536-entry
vocabulary. Our corpora are QUIC and TLS, where the payload is encrypted, and
our records retain sizes, directions and inter-arrival times but no bytes.
Feeding it header-derived features would not be ET-BERT; it would be a
differently-sized transformer carrying the name. Its published pre-training
corpus is also not redistributed. We therefore do not report a number for it,
rather than reporting one we could not obtain fairly. The same reasoning
applies to YaTC, CLE-TFE, MIETT, FlowletFormer and PacketCLIP; the artifact
documents each case individually.

This is a real limitation and it bounds our claims: we do not assert
state-of-the-art accuracy, and \S\ref{sec:results} does not contain the
strongest published models.

\subsection{Evaluation protocol}

The methodological question we raise is not new in machine learning generally.
Group-aware cross-validation is standard wherever observations cluster, and
tree-based models are known to outperform neural networks on tabular
features~\cite{grinsztajn2022trees} --- a result we reproduce in
\S\ref{sec:results} in a setting where it is rarely tested.

What is specific to encrypted traffic is \emph{which} grouping matters, and
that this differs by corpus in a way that is not visible from the label space.
A dataset of per-application packet captures and a dataset of backbone flows
can carry the same label taxonomy, the same feature schema and the same class
balance, and require entirely different partitions to be evaluated honestly.
We are not aware of prior work that measures this contrast directly on two
corpora, or that provides a pre-training diagnostic for choosing the axis.

\subsection{Open-set recognition}

Deployments encounter applications absent from training. Rejection by maximum
softmax probability is the usual baseline; energy scores~\cite{liu2020energy}
and Mahalanobis distance to class-conditional Gaussians are common
alternatives. We evaluate all of these on a single frozen embedding
(\S\ref{sec:results}), which isolates the scoring rule from the representation
and produces the negative result in \S\ref{sec:negative}: the rejection
advantage we measured is a property of the rule, not of the metric training we
added to produce it.
